\section{Results: Evidence Required}
\textbf{No empirical results are available in this manuscript version.} This section is an explicit reporting contract, not a table of observations.

\begin{table}[htbp]
\caption{Planned participant-level continuous valence evaluation; no observations are reported}
\centering
\begin{tabular}{llll}
\toprule
Corpus & Montage & CCC & MSE \\
\midrule
Eligible corpus & Full & Pending & Pending\\
Eligible corpus & Sparse $k=4$ & Pending & Pending\\
Held-out corpus & Compatible subset & Not eligible & Not eligible\\
\bottomrule
\end{tabular}
\end{table}

CCC is undefined rather than zero for participants whose observed valence ratings are constant. MSE must be reported for each eligible held-out participant against a training-only mean-response baseline. A categorical analysis (balanced accuracy and Brier score) would require a separate prespecified transformation and is not part of this continuous-outcome table.

\subsection{Prespecified reporting ledger}
For each dataset and montage, a future results release must report the
number of source participants, independent test participants, eligible
test trials, removed trials, electrode coordinates and reference scheme.
The reporting ledger will state whether participant-only, jointly
participant--stimulus-disjoint and held-out-corpus evaluations were
possible. All numbers remain pending until rights and input integrity
have been verified.

\begin{table}[htbp]
\caption{Prospective evidence gates; no gate is claimed to have passed}
\centering
\begin{tabular}{p{0.39\linewidth}p{0.48\linewidth}}
\toprule
Gate & Required record \\
\midrule
Rights and provenance & Version, licence, attribution, file checksums \\
Montage comparability & Coordinates, references, overlap, exclusions \\
Split integrity & Participant and stimulus IDs, fold-manifest hash \\
Train-only processing & Fit/transform audit, hyperparameter selection \\
Model comparisons & Constant, ridge, full and sparse, negative controls \\
Statistical reporting & CCC eligibility, paired MSE, interval method \\
\bottomrule
\end{tabular}
\end{table}

The reporting release must list undefined CCC values separately from
missing predictions and failed processing runs. A model may beat a
constant predictor while failing stimulus-only or participant-identity
controls; such a result cannot be described as transferable neural
decoding. Cross-corpus rows must not display numeric results if target
and source measurement scales or montages lack an auditable mapping.

The final report will contain participant-level paired differences and intervals, the exact number of legal datasets and unique participants, montage/reference exclusions, train-only preprocessing audit, loss of utility versus electrode count and a corpus-by-montage comparison only where scientifically identifiable. Negative controls and failure cases must be shown alongside positive results. Any claim unsupported by an actual data and code run will be omitted.
